% ---------------------------------------------------------------
% Corps du guide de mise en oeuvre de la phase 1.
% Partage par guide_phase1.tex (document autonome) et par
% cours_theorie.tex (annexe). Ne pas dupliquer : editer ici.
%
% Le document appelant doit fournir :
%   packages : enumitem, float, longtable, booktabs, graphicx, array,
%              textcomp
%   environnements : attention, critere
%   macros : \eps  (= \varepsilon_r)
%   graphicspath incluant figures/
% ---------------------------------------------------------------

%=======================================================================
\section{Avant de commencer}
%=======================================================================

\subsection{Ce que ce guide suppose}

Il suppose que vous savez : lire un schéma bloc, utiliser un multimètre,
installer un logiciel sous Windows, écrire ou modifier un script Python simple,
et visser un connecteur. Il ne suppose \emph{aucune} connaissance en
hyperfréquences. Chaque notion nécessaire est introduite là où elle sert.

\subsection{Sécurité}

Les puissances mises en jeu sont de l'ordre du milliwatt --- mille fois moins
qu'un téléphone portable. \textbf{Il n'y a aucun risque d'exposition.} Les deux
vrais risques sont matériels :

\begin{itemize}[nosep]
\item \textbf{Les décharges électrostatiques} détruisent les circuits
      hyperfréquences. Portez un bracelet antistatique relié à la terre, ou à
      défaut touchez une masse métallique avant chaque manipulation.
\item \textbf{Les alimentations inversées} détruisent instantanément un VCO.
      Vérifiez chaque tension au multimètre \emph{avant} de connecter la carte.
\end{itemize}

\subsection{Ne pas serrer les connecteurs à la main}

Un connecteur SMA se serre à $0{,}5$~N$\cdot$m, avec une clé dynamométrique.
Trop serré, il se déforme et la mesure dérive ; pas assez, le contact est
intermittent et \emph{tout} le reste devient incompréhensible. Une clé
dynamométrique SMA coûte une trentaine d'euros ; c'est le meilleur achat du
projet.

%=======================================================================
\section{Objectif et critères de réussite}
%=======================================================================

La phase 1 doit répondre à deux questions, et à elles seules. Tout ce qui n'y
contribue pas est hors périmètre.

\begin{critere}
\textbf{Q1 --- Linéarité.} La rampe balaie-t-elle $6{,}0$ à $8{,}0$~GHz en
$100\,\mu$s avec une non-linéarité inférieure à $0{,}1\,\%$ ?\\[1mm]
\textbf{Q2 --- Mesure complexe.} Peut-on en tirer un coefficient de réflexion
$\Gamma(f)$ \emph{calibré}, dont l'amplitude et la phase concordent avec celles
mesurées par le LiteVNA sur la même cible ?
\end{critere}

\subsection{D'où vient le seuil de $0{,}1\,\%$}

L'épaisseur se déduit d'un \emph{écart} de retards. Une erreur relative sur la
pente de la rampe se reporte donc directement en erreur relative sur
l'épaisseur.

\begin{center}
\includegraphics[width=\textwidth]{fig_p1_critere.png}
\end{center}

\begin{table}[H]
\centering
\begin{tabular}{lccc}
\toprule
Non-linéarité & $d = 1$~mm & $3$~mm & $5$~mm \\
\midrule
$1\,\%$ & $0{,}010$~mm & $0{,}030$~mm & $0{,}050$~mm \\
$\mathbf{0{,}1\,\%}$ & $0{,}001$~mm & $0{,}003$~mm & $\mathbf{0{,}005}$~mm \\
$0{,}01\,\%$ & $0{,}0001$~mm & $0{,}0003$~mm & $0{,}0005$~mm \\
\bottomrule
\end{tabular}
\end{table}

À $0{,}1\,\%$, le biais introduit vaut $0{,}005$~mm sur une couche de $5$~mm,
soit vingt fois moins que l'objectif de $0{,}1$~mm. C'est confortable, et c'est
ce qu'atteint couramment une boucle asservie.

\subsection{Ce que la phase 1 ne fait pas}

Elle ne cherche ni la compacité, ni la vitesse de balayage, ni l'intégration
mécanique, ni la conformité réglementaire. Le montage sera encombrant, alimenté
par plusieurs blocs de laboratoire, et posé sur une table. C'est voulu.

%=======================================================================
\section{Principe : pourquoi ce montage, et pourquoi il est simple}
%=======================================================================

\subsection{Le principe en une phrase}

On émet un signal dont la fréquence monte linéairement, on le mélange avec
l'écho qui revient, et la \emph{différence} de fréquence obtenue est
proportionnelle à la distance de la cible. Un aller-retour de $150$~mm donne
$20$~kHz : une fréquence audible.

\begin{attention}
C'est le point qui rend cette phase abordable : bien que le signal émis soit à
$8$~GHz, \textbf{tout ce que l'on mesure est de l'audio}. Un ordinateur muni
d'une carte son convient. Aucun oscilloscope rapide, aucun analyseur de spectre
hyperfréquence n'est indispensable.
\end{attention}

\subsection{Pas besoin de démodulateur IQ}

Un radar de ce type est souvent décrit avec un démodulateur en quadrature, qui
livre deux voies dites I et Q. C'est nécessaire lorsqu'il faut distinguer une
cible en avant d'une cible en arrière. Ici, toutes les cibles sont à des
distances positives : \textbf{un simple mélangeur suffit}, et la phase de la
réflexion --- dont nous avons besoin --- est intégralement préservée dans le
signal audio. Cela économise le composant le plus cher et le plus difficile à
trouver au-dessus de $6$~GHz.

\subsection{Architecture}

\begin{center}
\includegraphics[width=\textwidth]{fig_p1_banc.png}
\end{center}

Le signal parcourt le chemin suivant. La \textbf{boucle à rampe} commande la
tension de réglage d'un \textbf{oscillateur} qui produit le balayage
$6$--$8$~GHz. Un \textbf{diviseur} en prélève une part qui servira de
référence ; le reste est amplifié et émis. L'écho reçu est amplifié par un
\textbf{amplificateur faible bruit} puis \textbf{mélangé} à la référence. Il en
sort un signal audio, filtré puis numérisé.

Les repères TP1 à TP5 sont des points de test : chacun s'obtient en dévissant
un câble SMA et en y branchant l'instrument de mesure.

%=======================================================================
\section{Nomenclature}
%=======================================================================

\begin{attention}
Les références ci-dessous sont des \emph{candidats}, cohérents avec le besoin au
moment de la rédaction. Les catalogues évoluent, et certaines cartes
d'évaluation sont fabriquées en petites séries. \textbf{L'étape~0 du projet est
de vérifier disponibilité, référence exacte de commande et prix} sur
\texttt{mouser.fr} avant tout achat. Ne commandez jamais un composant sans avoir
ouvert sa fiche technique et vérifié sa plage de fréquence.
\end{attention}

\begin{longtable}{>{\raggedright\arraybackslash}p{2.6cm}%
                  >{\raggedright\arraybackslash}p{5.4cm}%
                  >{\raggedright\arraybackslash}p{4.0cm}c}
\toprule
Fonction & Exigence & Candidat & Ordre de prix \\
\midrule
\endfirsthead
\toprule
Fonction & Exigence & Candidat & Ordre de prix \\
\midrule
\endhead
Boucle à rampe & synthétiseur fractionnaire jusqu'à $13$~GHz, générateur de
  rampe matériel & \textbf{EVAL-ADF4159} (version à VCO externe) & $300$--$500$~\texteuro \\[1mm]
Oscillateur & couvre $6{,}0$--$8{,}0$~GHz, sortie SMA, bruit de phase
  $< -95$~dBc/Hz à $100$~kHz & carte d'évaluation \textbf{HMC586LC4B}
  ($4$--$8$~GHz) & $200$--$400$~\texteuro \\[1mm]
Diviseur & $2$ voies, $0{,}5$--$8{,}5$~GHz, SMA & \textbf{ZX10-2-852+}
  (Mini-Circuits) & $60$--$100$~\texteuro \\[1mm]
Amplificateur d'émission & couvre $6$--$8$~GHz, gain $\sim 15$~dB, SMA
  & \textbf{ZX60-83W+} ($50$~MHz--$8$~GHz) & $100$--$150$~\texteuro \\[1mm]
Amplificateur faible bruit & couvre $6$--$8$~GHz, facteur de bruit $< 3$~dB
  & \emph{à choisir dans la même famille} & $100$--$200$~\texteuro \\[1mm]
Mélangeur & double équilibré, RF et OL couvrant $6$--$8$~GHz, FI continue à
  $1$~MHz & \emph{connectorisé SMA, niveau OL $+7$ à $+13$~dBm}
  & $100$--$200$~\texteuro \\[1mm]
Bande de base & filtre passe-bas $\sim 100$~kHz + amplificateur à gain réglable
  & module audio ou carte à monter & $50$--$150$~\texteuro \\[1mm]
Numériseur & $\ge 500$~kSa/s, $16$~bits, entrée différentielle
  & interface audio ou carte d'acquisition & $150$--$400$~\texteuro \\[1mm]
Antennes ($\times 2$) & couvrent $6$--$8$~GHz, faible encombrement
  & cornet ou antenne UWB à connecteur SMA & $100$--$300$~\texteuro \\[1mm]
Ligne à retard & $3$~m de câble coaxial faible perte, SMA
  & \emph{indispensable au test de linéarité} & $50$--$100$~\texteuro \\[1mm]
Câbles et accessoires & câbles SMA courts, adaptateurs, charges $50\,\Omega$,
  atténuateurs $3$ et $10$~dB & & $200$--$400$~\texteuro \\[1mm]
Alimentations & $2$ à $3$ voies de laboratoire, $0$--$15$~V, limitation en
  courant & & $200$--$400$~\texteuro \\[1mm]
Clé dynamométrique SMA & $0{,}5$~N$\cdot$m & & $30$--$60$~\texteuro \\
\midrule
\multicolumn{3}{l}{\textbf{Total indicatif}} & $\mathbf{1{,}6}$--$\mathbf{3{,}5}$~\textbf{k\texteuro} \\
\bottomrule
\end{longtable}

\subsection{Instruments de mesure}

\begin{table}[H]
\centering
\begin{tabular}{>{\raggedright\arraybackslash}p{4.6cm}%
                >{\raggedright\arraybackslash}p{9.4cm}}
\toprule
Instrument & Nécessité \\
\midrule
Multimètre & \textbf{indispensable} --- vérification des alimentations \\
Carte son ou acquisition & \textbf{indispensable} --- c'est l'instrument
  principal, le signal utile étant audio \\
LiteVNA (déjà possédé) & \textbf{très utile} jusqu'à $6{,}3$~GHz : vérification
  des câbles, du diviseur, des antennes, et référence de comparaison pour Q2 \\
Analyseur de spectre à $8$~GHz & \emph{souhaitable, non indispensable} ---
  le test de linéarité (§\ref{sec:testA}) s'en passe entièrement. À emprunter
  ou louer pour une journée si possible \\
\bottomrule
\end{tabular}
\end{table}

\begin{attention}
Le LiteVNA s'arrête à $6{,}3$~GHz et ne peut donc pas observer directement le
haut de la rampe. Ce n'est pas gênant : le test décisif convertit le balayage
en un signal audio, et se mesure sans aucun instrument hyperfréquence.
\end{attention}

%=======================================================================
\section{Sous-système A --- la source de rampe}
\label{sec:testA}
%=======================================================================

C'est le c\oe ur du projet et le seul véritable risque. Consacrez-lui la moitié
du temps.

\subsection{Rôle}

Produire un signal dont la fréquence monte de $6{,}0$ à $8{,}0$~GHz en
$100\,\mu$s, de façon aussi linéaire que possible, et recommencer $5\,000$ fois
par seconde.

\subsection{Pourquoi il faut une boucle asservie}

Un oscillateur commandé en tension possède une caractéristique
tension--fréquence \emph{fortement courbe} : appliquer une rampe de tension
donne une rampe de fréquence tordue, avec plusieurs pour-cent d'erreur. Les
réalisations amateur procèdent ainsi et se contentent de $100$~MHz d'excursion,
où la courbure reste tolérable. Sur $2$~GHz, c'est impossible.

La boucle à verrouillage de phase résout cela en \emph{asservissant} la
fréquence : elle la compare en permanence à une référence à quartz et corrige.
La linéarité ne dépend plus de l'oscillateur mais de la référence, stable à
quelques parties par million.

\subsection{Documentation à rassembler}

\begin{enumerate}[nosep]
\item La fiche technique de l'ADF4159 --- comprendre les registres de rampe
      (déviation, pas, durée).
\item Le guide utilisateur de la carte d'évaluation --- brochage,
      alimentations, connecteur de programmation.
\item Le logiciel \textbf{ADIsimPLL} du fabricant, gratuit. Il calcule le filtre
      de boucle et \emph{simule la rampe}.
\item La fiche technique du VCO --- plage de fréquence, sensibilité de réglage
      (MHz/V), plage de tension admissible.
\end{enumerate}

\subsection{Installation}

\begin{enumerate}
\item Alimentez la carte PLL seule, sans le VCO. Vérifiez chaque tension au
      multimètre.
\item Installez le logiciel de contrôle. Vérifiez la communication : un
      registre lu doit renvoyer ce qui y a été écrit.
\item \textbf{Concevez le filtre de boucle avec ADIsimPLL}, en saisissant la
      sensibilité réelle du VCO. Le logiciel fournit les valeurs des composants.
      C'est la seule étape qui demande de la méthode ; le logiciel fait le
      calcul.
\item Connectez le VCO. Vérifiez d'abord qu'il oscille en lui appliquant une
      tension de réglage fixe.
\end{enumerate}

\subsection{Test A1 --- le balayage existe-t-il ?}

\textbf{Montage :} sortie du VCO (TP2) vers un analyseur de spectre, en mode
mémoire de maximum. Si vous n'avez pas d'analyseur, passez directement au test
A2, qui est plus important.

\textbf{Attendu :} un plateau plat entre $6{,}0$ et $8{,}0$~GHz.

\begin{center}
\includegraphics[width=\textwidth]{fig_p1_tests.png}
\end{center}

\begin{critere}
Plateau présent sur toute la plage, plat à $\pm 3$~dB, sans trou ni raie
parasite plus de $20$~dB au-dessus du plateau.
\end{critere}

\subsection{Test A2 --- la linéarité, sans instrument hyperfréquence}

C'est \textbf{le test décisif de toute la phase 1}. Son élégance tient à ce
qu'il n'exige rien d'autre que le montage lui-même.

\textbf{Principe.} On envoie la rampe dans un long câble et on mélange la sortie
du câble avec la rampe d'origine. Le câble introduit un retard connu, donc une
fréquence de battement connue. \emph{Si la rampe est parfaitement linéaire, ce
battement est une raie pure.} Toute non-linéarité l'élargit. La largeur de la
raie mesure donc directement la linéarité --- et cette raie est à $280$~kHz,
c'est-à-dire dans la bande d'une carte son.

\textbf{Montage :}
\begin{enumerate}[nosep]
\item Sortie VCO $\to$ diviseur.
\item Voie 1 du diviseur $\to$ \textbf{3~m de câble} $\to$ entrée RF du
      mélangeur.
\item Voie 2 du diviseur $\to$ entrée OL du mélangeur.
\item Sortie FI du mélangeur $\to$ filtre $\to$ numériseur.
\end{enumerate}

Aucune antenne, aucun amplificateur : uniquement le diviseur, le câble et le
mélangeur.

\textbf{Traitement :} enregistrez $100$ rampes, appliquez à chacune une
transformée de Fourier, et mesurez la largeur de la raie à $-6$~dB.

\begin{critere}
La raie de battement se situe à $280 \pm 10$~kHz et sa largeur à $-6$~dB est
inférieure à $\mathbf{3}$~\textbf{kHz}, soit $1\,\%$ de sa fréquence.\\[1mm]
\emph{Justification :} pour un retard fixe, l'élargissement relatif de la raie
est égal à la non-linéarité relative de la rampe. Une raie large de $1\,\%$
correspond donc à $1\,\%$ de non-linéarité --- déjà acceptable pour
$0{,}03$~mm de biais. Viser $0{,}3$~kHz ($0{,}1\,\%$) pour la marge complète.
\end{critere}

Si la raie est large de plusieurs dizaines de kilohertz, la boucle n'asservit
pas correctement : reprenez le filtre de boucle dans ADIsimPLL, en réduisant la
bande passante de boucle.

%=======================================================================
\section{Sous-système B --- la voie d'émission}
%=======================================================================

\textbf{Rôle.} Prélever une part du signal pour la référence, amplifier le
reste et l'émettre.

\textbf{Test B1 --- le diviseur.} Mesurable au LiteVNA jusqu'à $6{,}3$~GHz :
$-3{,}5$~dB environ sur chaque voie, isolation meilleure que $18$~dB entre
voies. Un écart franc signale un connecteur mal serré.

\textbf{Test B2 --- l'amplificateur.} Alimentez-le, injectez un signal connu,
mesurez le gain. Vérifiez le courant consommé : conforme à la fiche technique à
$\pm 20\,\%$.

\begin{attention}
Ne branchez \textbf{jamais} un amplificateur sans charge à sa sortie : une
sortie ouverte peut le détruire. Une charge $50\,\Omega$ vissée est la position
de repos.
\end{attention}

%=======================================================================
\section{Sous-système C --- la voie de réception}
%=======================================================================

\textbf{Rôle.} Amplifier l'écho, puis le mélanger à la référence.

\textbf{Test C1 --- l'amplificateur faible bruit.} Même procédure que B2.

\textbf{Test C2 --- le mélangeur.} Injectez deux signaux de fréquences voisines
sur RF et OL, vérifiez que la sortie FI porte bien leur différence.

\begin{attention}
Un mélangeur exige un niveau de référence \emph{précis} sur son entrée OL
(typiquement $+7$ à $+13$~dBm). Trop faible, il ne fonctionne pas ; trop fort,
il sature. C'est le rôle des atténuateurs de la nomenclature : ajustez le
niveau avec un atténuateur de $3$ ou $10$~dB plutôt qu'en modifiant le montage.
\end{attention}

%=======================================================================
\section{Sous-système D --- bande de base et numérisation}
%=======================================================================

\textbf{Rôle.} Filtrer et amplifier le signal audio, puis le numériser en
synchronisme avec la rampe.

\textbf{La synchronisation est le point délicat.} Chaque enregistrement doit
commencer au début d'une rampe, sinon les rampes successives ne peuvent pas être
moyennées. La carte PLL fournit un signal indiquant le début de rampe :
utilisez-le comme déclenchement du numériseur.

\textbf{Test D1.} Injectez une sinusoïde connue de $20$~kHz à l'entrée. Vérifiez
que la transformée de Fourier la restitue à la bonne fréquence et à la bonne
amplitude, et que le plancher de bruit descend de $3$~dB chaque fois que l'on
double le nombre de moyennes --- preuve que le moyennage est cohérent.

%=======================================================================
\section{Intégration progressive}
%=======================================================================

N'assemblez jamais tout d'un coup. À chaque étape, un seul élément nouveau.

\begin{table}[H]
\centering
\begin{tabular}{clp{7.4cm}}
\toprule
Étape & Ajout & Vérification \\
\midrule
1 & PLL + VCO & test A1 et A2 \\
2 & + diviseur, câble de $3$~m & test A2 complet, raie à $280$~kHz \\
3 & + bande de base et numériseur & test D1, moyennage cohérent \\
4 & + amplificateur d'émission & niveau à TP3 conforme \\
5 & + LNA & gain total mesuré \\
6 & + antennes & premier écho réel \\
\bottomrule
\end{tabular}
\end{table}

%=======================================================================
\section{Plan de test système}
%=======================================================================

\subsection{Test S1 --- justesse de l'échelle des distances}

\textbf{Montage :} antennes face à une plaque métallique, à $100$, $200$ puis
$300$~mm mesurés au réglet.

\textbf{Attendu :} un pic unique dont la position croît linéairement avec la
distance réelle (panneau C de la figure des tests).

\begin{critere}
Écart entre distance mesurée et distance réelle inférieur à $2\,\%$, et
\emph{constant} en pourcentage. Un écart constant en pourcentage est une erreur
de pente de rampe, corrigible par le calcul ; un écart variable révèle une
non-linéarité résiduelle.
\end{critere}

\subsection{Test S2 --- calibration}

Appliquez la même procédure que celle déjà en service sur le LiteVNA
(\texttt{calibration.py}) : mesurez un court-circuit, un circuit ouvert et
une charge au plan de référence, résolvez les trois termes d'erreur, puis
vérifiez que la charge, une fois corrigée, tombe sous le plancher.

\begin{critere}
Après correction, la réflexion résiduelle de la charge est inférieure à
$-30$~dB.
\end{critere}

\subsection{Test S3 --- comparaison au LiteVNA (le test final)}

C'est celui qui répond à Q2. Choisissez une cible stable : d'abord une plaque
métallique à $150$~mm, puis une lame de \textbf{PMMA} (plexiglas,
$\eps = 2{,}6$) posée sur un bloc de polystyrène expansé, rien derrière. Une
lame libre dans l'air a deux interfaces de coefficients égaux au signe près :
ses deux échos ont donc la même amplitude, ce qui en fait la cible la plus
facile qui soit. Une épaisseur de $50$~mm convient ; mesurez-la au pied à
coulisse en cinq points, et demandez du PMMA \emph{coulé}, dont la tolérance
d'épaisseur est bien meilleure que celle de l'extrudé.

\begin{enumerate}[nosep]
\item Mesurez $\Gamma(f)$ avec le LiteVNA calibré, sur $6{,}0$--$6{,}3$~GHz
      (le haut de sa plage).
\item Mesurez la même cible avec le banc FMCW, calibré, sur la même
      sous-bande.
\item Superposez amplitude et phase.
\end{enumerate}

\begin{critere}
Sur la sous-bande commune $6{,}0$--$6{,}3$~GHz, l'écart entre les deux mesures
est inférieur à $\mathbf{1}$~\textbf{dB} en amplitude et $\mathbf{5^\circ}$ en
phase, après retrait d'un retard et d'un gain complexe constants.
\end{critere}

\begin{attention}
Ce test est le plus important du guide. Il ne vérifie pas une performance : il
vérifie que le banc mesure \emph{la même grandeur physique} qu'un instrument de
métrologie déjà calibré. Tant qu'il n'est pas passé, aucun résultat du banc
n'est exploitable.
\end{attention}

%=======================================================================
\section{Décision de fin de phase}
%=======================================================================

\begin{table}[H]
\centering
\begin{tabular}{>{\raggedright\arraybackslash}p{5.0cm}%
                >{\raggedright\arraybackslash}p{9.4cm}}
\toprule
Résultat & Conclusion \\
\midrule
A2 $< 0{,}3$~kHz, S3 conforme & \textbf{Risque levé.} Engager la phase 2
  (prototype intégré) \\[1mm]
A2 entre $0{,}3$ et $3$~kHz, S3 conforme & \textbf{Risque levé avec réserve.}
  Utilisable ; améliorer la boucle en phase 2 \\[1mm]
A2 $> 3$~kHz malgré plusieurs filtres de boucle & \textbf{Reconsidérer.}
  Essayer une excursion réduite ($6$--$7$~GHz), ou une architecture à rampe
  numérique \\[1mm]
S3 non conforme, A2 conforme & Le problème est dans la réception ou la
  calibration, non dans la rampe. Reprendre les tests C et S2 \\
\bottomrule
\end{tabular}
\end{table}

\begin{attention}
Une phase 1 qui conclut « non » en trois mois et quelques milliers d'euros est
un \emph{succès} : elle évite une phase 2 d'un an. Le seul mauvais résultat est
un résultat ambigu, faute d'avoir suivi les critères chiffrés.
\end{attention}

%=======================================================================
\section{Pièges les plus fréquents}
%=======================================================================

\begin{enumerate}
\item \textbf{Confondre plage de l'oscillateur et excursion balayée.} Un VCO
      « $4$--$8$~GHz » ne garantit pas que la boucle balaie $2$~GHz d'un trait.
      C'est précisément ce que teste A2.
\item \textbf{Câble déplacé entre deux mesures.} Toute flexion change la phase
      et ruine la calibration. Fixez les câbles avec du ruban adhésif.
\item \textbf{Niveau OL incorrect.} Cause la plus fréquente d'un mélangeur qui
      « ne marche pas ».
\item \textbf{Absence de déclenchement.} Sans signal de début de rampe, le
      moyennage détruit le signal au lieu de l'améliorer. Symptôme : le niveau
      \emph{baisse} quand on moyenne.
\item \textbf{Antennes trop proches.} À moins de $80$~mm, l'écho d'antenne et
      celui de la cible se confondent. Travaillez d'abord à $150$--$300$~mm ;
      la mesure à $20$~mm relève de la phase 2.
\item \textbf{Vouloir tout mesurer d'un coup.} Le tableau d'intégration
      progressive existe pour cela : un seul élément nouveau à la fois.
\end{enumerate}
